% !TEX root = ./main.tex
\documentclass[a4paper]{article}
% Course language: use [english] or [german].
\usepackage[german]{ajnotes}
\title{\textbf{Financial Accounting — Warenkonten}}
\begin{document}
    \maketitle
    % Attach the given problem sheet: drop it in this folder as sheet.pdf and
    % un-comment the next line (pdfpages is already loaded by the preamble).
    % \includepdf[pages=-]{sheet.pdf}
    Warenlager konto ist immer korrekt zu einstandspreisen. Nach verkauf werden Warenaufwände
    Immer Werte und nicht Stückzahlen erfasst.

    \exercise{14.1}
    \subexercise{a}
    \begin{itemize}
        \item  Delta lager = Endbestand - Anfangsbestand oder Endbestand = Anfangsbestand + delta Lager
        \item Wareneinkauf = Warenaufwand + delta Lager
        \item Warenertrag = Warenaufwand + Bruttogewinn
    \end{itemize}
        \subexercise{b}
  \begin{itemize}
      \item AB = EB - Delta Lager = 540k
      \item Wareneinkauf = Warenaufwand + delta Lager
      \item Warenertrag = Warenaufwand + Bruttogewinn
  \end{itemize}
    \subexercise{c}
\begin{itemize}
      \item EB = Wareneinkauf 
      \item 
      \item 
  \end{itemize}
    \exercise{14.2}
    
    \exercise{14.3}
    \subexercise{a}

\end{document}
